\chapter{Literature Review}\label{ch:literature_review}

This chapter positions the thesis around the decisions that must be made by a
robust cell-free ISAC transmitter. Rather than treating ISAC architectures,
sensing metrics, and optimization methods as separate topics, the review
follows the same chain as the proposed design: cooperation creates useful
spatial degrees of freedom, QoS protection constrains their use, target-state
accuracy determines the sensing requirement, and topology recovery determines
whether the computed solution is physically implementable.

\section{From Coexistence to Cooperative ISAC}

ISAC emerged from the observation that wireless communication and sensing rely
on the same electromagnetic propagation, spectrum, and radio hardware. Early
single-base-station systems therefore studied waveform sharing, communication--
sensing trade-offs, and self-interference management
\autocite{liuSurveyFundamentalLimits2022,huaOptimalTransmitBeamforming2023,heIntegratedSensingFullduplex2023}.
\autocite{liuIntegratedSensingCommunications2022,liuJointRadarCommunication2020}
These studies establish the core trade-off, but a single transmitter provides
limited geometric diversity and can be vulnerable to cell-edge blockage.

Cooperative and multistatic ISAC extend the sensing and communication aperture
across multiple base stations or APs. Spatially distributed observations can
improve coverage, localization geometry, and interference management, provided
that the transmitters are coordinated
\autocite{cuiReviewMultibaseStation2025,mengCooperativeISACNetworks2025,liaoCooperativeBaseStation2025}.
\autocite{chengNetworkedBeamforming2024,maoCommunicationSensingRegion2024,liNetworkedISAC2024}
Cell-free massive MIMO takes this idea further: a large set of distributed APs
is centrally coordinated without fixed cell boundaries
\autocite{femeniasScalableCellfreeMassive2025,demirhanCellfreeISACMIMO2025}.
\autocite{maoCommunicationSensingRegion2024,renSecureCellfreeISAC2024}
The same cooperation that improves service flexibility also couples every
AP's communication and sensing power decisions. This is the architectural
setting of the present thesis.

\section{Design Requirements for Robust Multi-Target Sensing}

\subsection{Communication QoS under Imperfect CSI}

Communication QoS is commonly expressed through rate or SINR constraints.
Many ISAC formulations assume perfect CSI to simplify joint beamforming, but
estimation error, feedback delay, and quantization make this assumption fragile
in practice \autocite{yuLocationSensingBeamforming2022,chuIntegratedSensingCommunication2023}.
\autocite{heMSETrainingISAC2024,gongDigitalTwinCSI2025}
For a cell-free transmitter, channel uncertainty is particularly consequential:
the dedicated sensing covariances and multi-user communication beams share the
same AP power budgets and jointly determine the user interference level.
This motivates a norm-bounded worst-case SINR requirement rather than a
nominal-channel guarantee.

\subsection{Target-Specific Detection and Tracking Accuracy}

Sensing designs are often evaluated with sensing SINR and Cram\'er--Rao-bound
(CRB) metrics. Sensing SINR is a detection-related proxy, whereas the CRB
measures a lower bound on estimation error and is consequently useful for
localization and tracking-oriented design
\autocite{liaoCooperativeBaseStation2025,gengJointBeamformingCRBconstrained2025}.
\autocite{liuCRBBeamforming2022}
For multiple targets, however, the sensing resource assigned to one target
should not be credited to another merely because the network uses a shared
transmit infrastructure. The present model therefore assigns a known sensing
covariance to each target and forms its FIM and matrix PCRB from that covariance
alone. This choice makes the association and tracking requirements refer to
the same physical sensing resource.

\subsection{Sensing Topology and Physical Recovery}

Cell-free systems can authorize subsets of APs to allocate a sensing waveform
for a target, trading sensing geometry against finite AP power and implementation
cost. This turns sensing-topology authorization into a mixed-integer design variable rather
than a post-processing detail. Semidefinite relaxation is valuable for handling
beamforming structure, but a relaxed covariance or a continuous association
weight does not by itself establish an implementable design. A complete
workflow must recover rank-one beamformers and binary associations, then test
the recovered variables against the original communication, sensing, tracking,
and power constraints.

Recent work also broadens ISAC across waveform, hardware, and surface-assisted
architectures: OTFS exploits delay--Doppler structure for high-mobility sensing
and communication \autocite{yuanOTFSISAC2021}; one-bit transceiver design
addresses low-resolution mmWave hardware \autocite{linOneBitISAC2025}; and
STAR-RIS and holographic surfaces provide relevant near-field extensions
\autocite{wangSTARRISNearfield2025,zhangHolographicISAC2022}. At the system
layer, sensing-aware status updating offers a complementary scheduling view
\autocite{pengSensingCommunicationStatus2023}. These directions motivate the
scope of this thesis but are not assumed by its conventional multi-antenna AP
signal model.

\section{Research Gap and Thesis Position}

The preceding literature motivates a precise gap at the intersection of the
four requirements. Cooperative ISAC provides the relevant spatial degrees of
freedom; robust QoS prevents their allocation from relying on nominal CSI;
target-specific FIM/PCRB modelling makes dedicated sensing participation
meaningful; and recovery-aware optimization separates a feasible relaxed point
from a validated physical design. This thesis combines these elements in one
power-minimizing cell-free ISAC formulation. Chapter~\ref{ch:system_model}
defines the required resource separation, and Chapter~\ref{ch:solution_method}
develops the corresponding robust and recovery-aware optimization method.
